\chapter{Literature Survey}

\section{Overview}

This chapter reviews the existing work that is relevant to this project. The main areas are
sentiment analysis and emotion detection, handling ambiguity and uncertainty in natural language
processing, topological data analysis for text, and the specific problems that have been observed
with existing models on the GoEmotions dataset.

\section{Sentiment Analysis}

Sentiment analysis is the task of identifying whether a piece of text expresses a positive,
negative, or neutral opinion. It is one of the oldest and most studied problems in natural language
processing.

to individual words and then combining those scores to get an overall sentiment. This
is a well-known approach for early systems. These were designed specifically for social media
text and handled punctuation, capitalization, and emoticons as signals of sentiment intensity.
Such systems are fast and work reasonably well for clear, unambiguous text. However, they
often fail on sarcasm and do not have any mechanism to say when it is unsure.

Machine learning approaches improved on this. Systems based on Naive Bayes, Support Vector
Machines (SVM), and logistic regression used bag-of-words representations. These models
could learn from examples and generalize beyond a fixed vocabulary. They performed better on
benchmark datasets but still struggled with informal language, negations, and short posts where
context is limited.

The deep learning era brought significant improvements. Previous models applied convolutional
filters over word embeddings to capture local phrases and patterns. Recurrent networks,
especially LSTMs, captured longer dependencies in a sequence. Attention mechanisms allowed
models to focus on the most relevant parts of the input.

BERT~\cite{bert} changed the field significantly. It was pre-trained on a massive amount of text using
a masked language modeling objective, which gave it a deep understanding of context.
Fine-tuned BERT models reached state-of-the-art results on almost every sentiment and emotion
benchmark. RoBERTa and DeBERTa further refined this approach.

The problem is that all of these systems, from VADER to BERT fine-tuned models, \textbf{always
return a label}. They have no built-in way to say that the input is unclear. They are trained
to minimize classification error on a fixed label set, and that is all they do.

\section{Emotion Detection and the GoEmotions Dataset}

Emotion detection is a more specific version of sentiment analysis. The goal is to identify
which particular emotion a text expresses, such as joy, sadness, anger, fear, or disgust. This
is harder than binary sentiment classification because the boundaries between emotions are not
sharp.

The GoEmotions dataset~\cite{goemotions} is one of the most important resources for this task. It was
released by Google and contains around 58,000 Reddit comments labelled by crowd workers
with 28 emotion categories. Each comment can have multiple labels. The dataset is publicly
available and has been widely used in research.

Several important problems have been identified with models trained on GoEmotions:

The first is \textbf{class imbalance}. The Neutral class makes up a large portion of the dataset.
Many specific emotion categories like grief, nervousness, and pride have very few examples.
Models trained on this data tend to predict the majority classes much more often and perform
poorly on minority classes. This is visible in the classification reports of most published systems
on this dataset.

The second problem is \textbf{label ambiguity}. Because the dataset was labelled by crowd
workers, there is significant disagreement between annotators. A comment like ``Well, that
happened'' received different labels from different people. Some saw it as neutral, some as
surprised, some as slightly negative. The final label is determined by majority vote, which hides
this uncertainty. Models trained on such data are essentially learning to predict the most common
human interpretation, not the correct one, and there may not be a single correct one.

The third problem is the presence of \textbf{informal language}. Reddit comments contain slang,
abbreviations, and expressions that are not well-represented in standard training data. Words
like ``yolo'', ``fomo'', ``lol'', and ``idk'' carry emotional meaning but are often handled poorly
by models trained on more formal text. The embedding models used to represent these words
may not give them accurate positions in the semantic space.

Several BERT-based systems have been published on GoEmotions~\cite{bert_emotion, bert_model}. Most report accuracy in the
range of 0.46 to 0.58 on the full 28-class problem~\cite{goemotion_models, llm_emotion}. On the simplified 3-class version (Positive,
Negative, Neutral), results are higher but still reflect the difficulty of the dataset~\cite{goemotions}. Several studies have explored emotion detection using the GoEmotions dataset. These approaches primarily rely on transformer-based architectures and deep learning models to achieve strong classification performance.

\begin{table}[H]
\centering
\caption{Comparison of Existing Methods on the GoEmotions Dataset}
\label{tab:goemotions_comparison}
\begin{tabular}{|l|l|l|l|}
\hline
\textbf{Study} & \textbf{Method} & \textbf{Metric} & \textbf{Score} \\ \hline
Demszky et al. (2020) & BERT (baseline) & Macro-F1 & 0.46 \\ \hline
Sentence-BERT (2019) & Transformer Embedding & Accuracy & $\sim$82\% \\ \hline
BiLSTM Model (2021) & Deep Learning & Macro-F1 & $\sim$0.45 \\ \hline
Transformer Model (2021) & Fine-tuned Transformer & Micro-F1 & $\sim$0.72 \\ \hline
LLM-based Model (2023) & Transformer + Augmentation & Accuracy & $\sim$85\% \\ \hline
ModernBERT (2024) & Advanced Transformer & Accuracy & $\sim$86--87\% \\ \hline
\end{tabular}
\end{table}

From \textbf{Table~\ref{tab:goemotions_comparison}}, it can be observed that transformer-based models achieve high performance on the GoEmotions dataset. However, these approaches are primarily designed for classification and do not explicitly account for ambiguity in emotional expressions.

\section{Ambiguity and Uncertainty in Natural Language Processing}

The question of how to handle uncertainty in NLP predictions has been studied in several
different ways.

Bayesian deep learning approaches model uncertainty by treating model weights as
distributions rather than fixed values. At test time, dropout is kept active and multiple predictions
are made. The variance across these predictions estimates model uncertainty. This works but
requires multiple forward passes and adds computational cost.

Conformal prediction is another framework that produces prediction sets rather than single
labels. Instead of saying ``this is Positive'', a conformal predictor might say ``this is Positive or
Ambiguous, and I am 90\% confident the true label is in this set.'' This is more honest but
harder to explain to a non-technical user.

Selective prediction allows a model to abstain from making a prediction when it is below
a confidence threshold. This is close to what this project does, but it relies only on the
classifier's own confidence. It does not use any external signal about the structure of the data.

The key limitation of all these approaches is that they use the model's own output to
estimate uncertainty. If the model is confidently wrong, none of these methods will catch it.
Recent research on topological quantification of ambiguity~\cite{tda_ambiguity, blowfish} suggests
that geometric signals can provide a more robust measure of uncertainty. Other works such as
MARCH~\cite{march} have evaluated the intersection of ambiguity interpretation and
multi-hop inference, further highlighting the complexity of handling unclear inputs.
What is needed is an external signal that looks at the data itself, not just at what the model says.

\section{Topological Data Analysis in Natural Language Processing}

Topological Data Analysis (TDA) is a field of mathematics that studies the shape of data. The
main tool is persistent homology~\cite{tda_foundation}, which tracks the birth and death of topological features as
a parameter is varied. The key features are connected components (H0 features) and loops or
holes (H1 features).

The Vietoris-Rips complex is the most common construction used in practice. Given a set of
points and a scale parameter, it builds a simplicial complex by connecting points that are closer
than the scale. As the scale increases, new components connect and new loops appear and then
fill in. Persistent homology records which features appear and how long they persist. Features
that persist for a long time are considered genuinely structural, while features that appear and
disappear quickly are considered noise.

TDA has been applied to NLP in several ways. Zhu~\cite{tda_nlp} applied persistent homology to text
classification and showed that topological features captured information that was not present in
bag-of-words representations. Michel et al.~\cite{embedding_geometry} analyzed the topology of sentence embedding
spaces and found that the topology changes depending on the semantic content of the text. Recent
reviews~\cite{tda_review} have highlighted the potential of topological deep learning beyond
standard persistence. In affective computing, AffectiveTDA~\cite{affective_tda} has been used
to improve both analysis and explainability. Beyond NLP, TDA has found success in analyzing
complex signals in fields such as cardiology~\cite{tda_cardio} and neurophysiology, where it is
used as a tool for EEG processing~\cite{tda_eeg}. Tools such as giotto-tda~\cite{giotto},
teaspoon~\cite{teaspoon}, and TDAvec~\cite{tdavec} have made these complex mathematical
approaches more accessible to the machine learning community.

What this project does differently is apply persistent homology at the \textbf{neighbourhood
level at inference time}. Instead of computing topology over the whole dataset, it computes a
fresh persistence diagram for the 24 nearest neighbours of each new input. This local topology
captures whether the input is sitting in a clean region of the space or in a boundary region
between clusters. The number of H1 loops in this local diagram is used as a signal of ambiguity.

This approach has not been used for ambiguity detection in social media text, to the best
knowledge of the author. It is the main technical contribution of this project.

\section{Sentence Embeddings}

The system uses SentenceTransformers~\cite{sentence_bert} to convert text into dense vector representations.
The specific model used is \texttt{all-mpnet-base-v2}, which produces 768-dimensional vectors.
This model was trained to place semantically similar sentences close to each other in the vector
space.

Modern embedding models like this one are trained on very large datasets and learn to handle
most of the variation in natural language automatically. They do not require explicit
preprocessing steps like tokenization, stopword removal, or stemming to be performed manually.
In fact, applying such preprocessing can hurt performance because it removes information that
the embedding model would have used.

This is an important design choice in this project and is discussed in more detail in the
implementation chapter.

\section{Gaps in Existing Work and How This Project Addresses Them}

Despite the strong performance of transformer-based models on the GoEmotions dataset, existing approaches lack the ability to explicitly detect ambiguity in emotional expression. Most models are designed to assign a fixed label, even when the input contains mixed, unclear, or conflicting emotional signals.

This limitation is significant because real-world social media text often contains subtle, implicit, or ambiguous emotions that cannot be accurately represented by a single label.

Furthermore, while embedding-based methods capture semantic similarity, they do not analyze the structural complexity of the data distribution in the embedding space.

This project addresses this gap by integrating semantic embeddings with topological features to detect structurally complex regions in the embedding space. Instead of forcing a classification, the system identifies ambiguous inputs and assigns an ``Ambiguous'' label when necessary.
